 \documentclass[3p]{elsarticle}
%\documentclass[preprint]{elsarticle}

\usepackage{amsmath, amsfonts, amsthm, bm}
\usepackage{amssymb}

\theoremstyle{remark}
\newtheorem*{remark}{Remark}

\usepackage{graphicx}
\usepackage{svg}         

\usepackage{array}         
\usepackage{booktabs}      % \toprule \midrule \bottomrule
\usepackage{multirow}
\usepackage{threeparttable}

\usepackage{stfloats}

\usepackage{algorithm}
\usepackage{algorithmic}   

\usepackage[T1]{fontenc}
\usepackage{textcomp}
\usepackage{pifont}      
\usepackage{enumitem}

\usepackage{url}
\usepackage{pdfpages}
\usepackage{verbatim}
\usepackage{subcaption}    
\usepackage{multicol}
\usepackage{xcolor}        
\usepackage{soul}          
\usepackage{tcolorbox}

\usepackage{tabularx}           
\newcolumntype{Y}{>{\raggedright\arraybackslash}X}   
\newcolumntype{C}{>{\centering\arraybackslash}X}     

% ===========================
% ===========================
\newtheorem{theorem}{Theorem}
\newtheorem{definition}{Definition}
\newtheorem{lemma}{Lemma}
\newtheorem{corollary}{Corollary}
\newtheorem{assumption}{Assumption}

% ===========================
% ===========================

\journal{Computers and Electrical Engineering}
% ========================================================================
\usepackage{CJKutf8}

\biboptions{sort&compress}

\begin{document}

\section*{Summary of Revisions}

In accordance with the editor’s decision and the reviewers’ comments, we have carefully revised the manuscript. 
Below we summarize how each major comment has been addressed. 
This section is included only in the revised submission and will be removed in the final camera-ready version.

\textbf{Reviewer 1}
\begin{itemize}

    \item \textbf{Comment 1:} 
    The abstract is too long and should be tightened to emphasize the key contributions and quantitative findings.

    \textbf{Response:}
    We thank the reviewer for the helpful suggestion. The abstract has been thoroughly revised and significantly shortened. The updated version highlights CleVer’s core contributions and main results in a more concise and focused manner.

    \item \textbf{Comment 2:}
    The mathematical symbols (e.g., $F(t)$, $f(s)$) are reused with slightly varying meanings. A notation summary table should be added.

    \textbf{Response:}
    We appreciate the reviewer’s observation. 
    We have added a complete notation summary table (Table 2) near the beginning of the paper. 
    The symbols are now consistently defined, with clarified meanings and units where appropriate.

    \item \textbf{Comment 3:}
    Typographical issues: “we describe the system model” → 
    “Section~4 describes…”. The expression “V er(·)” should be consistently formatted.

    \textbf{Response:}
    Thank you for pointing this out. 
    We have revised the affected sentences to use correct section references. All occurrences of the verification operator have been standardized to $Ver(\cdot)$ for consistency throughout the manuscript.

    \item \textbf{Comment 4:}
    Several tables lack proper captions or units.

    \textbf{Response:}
    We have updated all tables to include complete captions, units, 
    and explanatory footnotes where necessary, and the tables now include explicit definitions of reported metrics and their measurement units.

    \item \textbf{Comment 5:}
    Figures 3 and 4 are informative but visually dense. 
    Consider simplifying flow arrows or splitting into subfigures.

    \textbf{Response:}
    We agree and have simplified Figures~3 and~4 to improve readability. 
    The revised versions reduce visual density by streamlining flow arrows and restructuring elements to make the diagrams easier to interpret.

    \item \textbf{Comment 6:}
    The conclusion requires a brief discussion of limitations, such as 
    dependence on blockchain reliability and assumptions about honest majority.

    \textbf{Response:}
    We appreciate this suggestion. 
    We have expanded the conclusion to include a concise discussion of 
    CleVer’s limitations, including its reliance on blockchain liveness 
    and stability, and the increased hardware demands for verifying 
    computation-intensive tasks. We also outline possible directions for future improvement.

\end{itemize}


\textbf{Reviewer 2}
\begin{itemize}

    \item \textbf{Comment 1:} 
    The manuscript contains grammatical and formatting issues. Specifically:
    (i) In Section 2.3, the subject–verb agreement in “None of these directions fully satisfy…” is incorrect.  
    (ii) Section 8 lacks a glossary table explaining symbolic variables.  
    (iii) When comparing to prior work, references should be re-cited instead of mentioning scheme names directly.  
    Several other grammatical issues remain.

    \textbf{Response:} 
    We thank the reviewer for the careful reading.  
    We have thoroughly revised the manuscript to correct all grammatical and formatting issues.  
    In Section 2.3, the verb has been corrected to “fully satisfies” to maintain proper subject–verb agreement.  
    Furthermore, a complete glossary of symbolic variables used in the experimental section has been added, including detailed definitions and units.  
    In Section 8, all prior schemes mentioned in the comparison table and explanatory text now include explicit citations.  
    Additional formatting issues (notation consistency, capitalization, and table/figure references) have also been corrected throughout the manuscript.
    
    \item \textbf{Comment 2:} 
    Only 45\% of references fall within 2021–2025, reducing the timeliness and relevance of the literature review.

    \textbf{Response:} 
    We have significantly updated the reference list to incorporate recent works from 2021–2025 in verifiable computation, optimistic rollups, anonymous verification, and staking-based incentive systems.  
    The proportion of recent citations has been substantially increased, strengthening the timeliness and relevance of the related-work section.

    \item \textbf{Comment 3:} 
    The Introduction suggests that not using ZKP is an advantage, but Section 6 employs ZKP for anonymous verification, creating an apparent contradiction.

    \textbf{Response:}
    We appreciate this valuable observation.  
    We have clarified both in the Introduction and in Section 6 that CleVer does \emph{not} rely on ZKP for computation verification.  
    The ZKP module is used solely as an auxiliary identity-protection component in the anonymous verifier staking and election protocol.  
    This usage follows a separate design path and does not affect CleVer’s core verification logic or its avoidance of ZK-based computation proofs.  
    An explicit explanation has been added in the contribution summary and Section 6 to ensure conceptual consistency.

    \item \textbf{Comment 4:} 
    Section 7’s security analysis does not clearly correspond to the threat model in Section 4.3.

    \textbf{Response:} 
    We have fully revised Section 4.3 to present a clearer and more structured threat model with three explicit categories:  
    (i) verifier identity exposure and targeted attacks,  
    (ii) dishonest or lazy verification behavior, and  
    (iii) incentive failures and challenge-window abuse.  
    In Section 7, the security analysis has been reorganized or referenced to match these three categories directly:  
    system anonymity (Section 7.1), honest verification enforcement (Section 7.2), and incentive compatibility with bounded blocking (Section 7.3).  
    This creates a clear one-to-one mapping between threats and defenses, resolving the logical gap noted by the reviewer.

    \item \textbf{Comment 5:} 
    The comparison in Section 8 mixes theoretical and implementation-based results, making the evaluation less consistent.

    \textbf{Response:} 
    At the beginning of Section 8, we have added a methodological clarification explaining that prior schemes such as Truebit, Arbitrum, and state-channel approaches report only asymptotic or analytical performance in their original papers.  
    Therefore, we present analytical comparisons for structural properties such as redundancy and finality.  
    CleVer’s implementation-level performance is evaluated separately in the subsequent subsections.  
    These two dimensions—analytical comparison for structural behavior and empirical benchmarks for realized performance—are complementary and together form a complete and consistent evaluation methodology.

    \item \textbf{Comment 6:} 
    The evaluation lacks detailed comparisons with recent verifiable computation frameworks.

    \textbf{Response:} 
    We have expanded the discussion in Section 8 to include comparisons against recent verification frameworks where meaningful metrics are available. For frameworks without publicly available prototypes, we provide analytical comparisons following standard practice in the literature.  
    This enhanced comparison highlights CleVer’s structural advantages—reduced dispute rounds, avoidance of redundant recomputation, flexible staking-based finality—while empirical experiments demonstrate its practical performance on real workloads. \\

\end{itemize}



\textbf{Reviewer 3}
\begin{itemize}

    \item \textbf{Comment 1: Terminology consistency.}  
    Standardize abbreviations and technical terms.

    \textbf{Response:}  
    We have reviewed the manuscript and standardized all terminology and abbreviations.  

    \item \textbf{Comment 2: Language and structure.}  
    Improve grammatical accuracy and sentence flow.

    \textbf{Response:}  
    We have performed a careful revision of the manuscript to correct grammatical errors and improve sentence structure.  
    Examples highlighted by the reviewer have been corrected.     Additional sentences throughout the paper have been refined for clarity and readability.

    \item \textbf{Comment 3: Figures and tables.}  
    Ensure sequential referencing, unify caption formats, and remove duplicate captions.

    \textbf{Response:}  
    All figures and tables have been checked for sequential numbering and consistent referencing.  

    \item \textbf{Comment 4: Experimental setup clarification.}  
    Explain how the snapshot interval threshold is determined (fixed or dynamically tuned).

    \textbf{Response:}  
    We have added a clarification in Section~8 to state that the snapshot interval is determined by the accumulated instruction weight, and that instruction weights can be dynamically tuned according to blockchain node execution capabilities.  

    \item \textbf{Comment 5: Notation and variable consistency.}  
    Unify symbol formatting such as $T_w$ vs.\ $Tw$.

    \textbf{Response:}  
    We have unified mathematical notation across the entire manuscript.  The notation table (Table~\ref{Notations}) provides a clean and authoritative reference for all symbols.

    \item \textbf{Comment 6: Anonymity scope explanation.}  
    Clarify whether verifier anonymity remains unlinkable across multiple tasks.

    \textbf{Response:}  
    A dedicated clarification has been added to Section~6.  
    We now explicitly state that verifier identities remain unlinkable across different tasks because  each participation uses a fresh commitment and nullifier, preventing cross-task linkage.  This emphasizes the strong anonymity guarantees of the protocol.

    \item \textbf{Comment 7: Reference formatting.}  
    Merge citation ranges and ensure ascending numerical order.

    \textbf{Response:}  
    All references have been reorganized to follow ascending numerical order.  
    Citation ranges have been merged into compact forms.  

    \item \textbf{Comment 8: Formatting consistency.}  
    Unify capitalization of section titles and remove duplicated explanations.

    \textbf{Response:}  
    Section titles have been standardized to title case.  
    Redundant or repeated explanations—such as duplicated content on static versus dynamic data—have been removed.  
    Overall formatting consistency has been improved across the manuscript.

\end{itemize}


\end{document}
